\documentclass[11pt,a4paper]{article}

% === Packages ===
\usepackage[margin=2.5cm]{geometry}
\usepackage{amsmath,amssymb,amsthm}
\usepackage{bm}
\usepackage{booktabs}
\usepackage{array}
\usepackage{hyperref}
\usepackage{enumitem}
\usepackage{xcolor}

% === Theorem environments ===
\newtheorem{definition}{Definition}[section]
\newtheorem{proposition}{Proposition}[section]
\newtheorem{corollary}{Corollary}[section]
\newtheorem{remark}{Remark}[section]
\newtheorem{theorem}{Theorem}[section]

% === Shortcuts ===
\newcommand{\vect}[1]{\bm{#1}}
\newcommand{\mat}[1]{\bm{#1}}
\newcommand{\feq}{f^{\mathrm{eq}}}
\newcommand{\fneq}{f^{\mathrm{neq}}}
\newcommand{\meq}{m^{\mathrm{eq}}}
\newcommand{\mneq}{m^{\mathrm{neq}}}
\newcommand{\Tmat}{\mat{T}}
\newcommand{\Mmat}{\mat{M}}
\newcommand{\Smat}{\mat{S}}
\newcommand{\Imat}{\mat{I}}
\newcommand{\tu}{\tilde{u}}
\newcommand{\dd}{\mathrm{d}}
\newcommand{\Real}{\mathbb{R}}

\title{Construction of the Central Moment Shift Operator $\Tmat(\tilde{\vect{u}})$\\
for D3Q19 MRT Lattice Boltzmann Methods:\\
Lattice-Reduced Polynomial Approach}

\author{Technical Report --- GILBM Project}
\date{March 2026}

\begin{document}
\maketitle

\begin{abstract}
We present a complete, self-contained derivation of the shift operator $\Tmat(\tilde{\vect{u}})$ that transforms raw moments to central moments in the D3Q19 Multiple-Relaxation-Time (MRT) Lattice Boltzmann framework.
The shift operator is defined via Dubois et al.\ (2015) as $\Tmat(\tilde{\vect{u}}) = \Mmat(\tilde{\vect{u}})\,\Mmat^{-1}(0)$, where $\Mmat(\tilde{\vect{u}})$ evaluates the d'Humi\`eres polynomial basis at shifted lattice velocities $\vect{e}_i - \tilde{\vect{u}}$.
We show that the d'Humi\`eres basis polynomials for D3Q19 contain degree-4 terms whose continuous polynomial extension is \emph{not unique} on the lattice.
We identify two choices --- the ``Dubois continuous'' form and the ``lattice-reduced'' form --- and prove that only the latter yields an algebraically exact group structure $\Tmat(\vect{u})\,\Tmat(-\vect{u}) = \Imat$.
Both forms produce identical stress-mode central moments and hence the same Navier--Stokes viscosity (Corollary~2.1 of Dubois et al.).
Full closed-form expressions for the 19$\times$19 $\Tmat(\tilde{\vect{u}})$ and $\Tmat^{-1}(\tilde{\vect{u}})$ are given.
\end{abstract}

\tableofcontents
\newpage

% ============================================================
\section{Introduction and Motivation}
% ============================================================

The MRT-CM (Central Moment) collision operator for lattice Boltzmann methods relaxes the non-equilibrium distribution in a frame co-moving with the local fluid velocity, thereby improving Galilean invariance compared to standard MRT relaxation in the rest frame.
The collision step takes the form
\begin{equation}\label{eq:collision}
  f_i^* = f_i - \bigl(\Mmat^{-1}\,\Tmat^{-1}(\vect{u})\,\Smat\,\Tmat(\vect{u})\,\Mmat\bigr)_{ij}\,(f_j - \feq_j),
\end{equation}
where $\Mmat$ is the moment transformation matrix, $\Smat = \mathrm{diag}(s_0,\ldots,s_{18})$ is the diagonal relaxation matrix, and $\Tmat(\vect{u})$ is the \emph{shift operator} that converts raw moments to central moments.

Dubois et al.\ (2015) introduced a general framework: given a polynomial basis $\{P_0, \ldots, P_{q-1}\}$, the shifted moment matrix is $M_{kj}(\tilde{\vect{u}}) = P_k(\vect{v}_j - \tilde{\vect{u}})$, and the shift operator is
\begin{equation}\label{eq:T-def}
  \Tmat(\tilde{\vect{u}}) = \Mmat(\tilde{\vect{u}})\,\Mmat^{-1}(0).
\end{equation}

However, Dubois et al.\ only derived \emph{explicit} closed-form expressions for D2Q9.
\textbf{No explicit construction for D3Q19 was given in their paper.}
This report fills that gap: we derive the complete $19\times 19$ shift operator for D3Q19 using the d'Humi\`eres (2002) polynomial basis, and we identify a critical subtlety --- the polynomial extension ambiguity --- that must be resolved to obtain correct algebraic properties.

% ============================================================
\section{The D3Q19 Lattice and d'Humi\`eres Basis}
% ============================================================

\subsection{Lattice velocities}

The D3Q19 lattice has velocity vectors $\vect{e}_i = (e_{ix}, e_{iy}, e_{iz})$ with components in $\{-1, 0, 1\}$:
\begin{equation}\label{eq:lattice}
\begin{aligned}
  &\vect{e}_0 = (0,0,0), \\
  &\vect{e}_{1,2} = (\pm 1,0,0),\quad \vect{e}_{3,4} = (0,\pm 1,0),\quad \vect{e}_{5,6} = (0,0,\pm 1), \\
  &\vect{e}_{7\text{--}10} = (\pm 1,\pm 1,0),\quad \vect{e}_{11\text{--}14} = (\pm 1,0,\pm 1),\quad \vect{e}_{15\text{--}18} = (0,\pm 1,\pm 1).
\end{aligned}
\end{equation}

\subsection{The lattice reduction identity}\label{sec:lattice-id}

Since all velocity components satisfy $e_{i\alpha} \in \{-1,0,1\}$, we have the fundamental identities:
\begin{equation}\label{eq:lattice-reduction}
  \boxed{e_{i\alpha}^{2n} = e_{i\alpha}^2, \qquad e_{i\alpha}^{2n+1} = e_{i\alpha}, \qquad \forall\, n \geq 1.}
\end{equation}
This means that on the lattice, $x^4 = x^2$, $x^3 = x$, $x^6 = x^2$, etc.
Consequently, any polynomial evaluated on the lattice can be reduced to a polynomial involving only the \textbf{19 independent monomials}:
\begin{equation}\label{eq:19-monomials}
  \{1,\; x,\; y,\; z,\; x^2,\; y^2,\; z^2,\; xy,\; xz,\; yz,\; x^2 y,\; x^2 z,\; xy^2,\; y^2 z,\; xz^2,\; yz^2,\; x^2 y^2,\; x^2 z^2,\; y^2 z^2\}.
\end{equation}

\begin{remark}
These 19 monomials span $\Real[\vect{e}]|_{\text{D3Q19}}$, the space of all polynomial functions restricted to the 19 lattice points.
Any polynomial in $x,y,z$ evaluated on D3Q19 is equivalent to a unique linear combination of these 19 monomials.
\end{remark}

\subsection{The d'Humi\`eres polynomial basis}

d'Humi\`eres et al.\ (2002) defined the moment matrix $\Mmat$ by its action on lattice velocities:
$M_{k,i} = P_k(e_{ix}, e_{iy}, e_{iz})$, where $\{P_k\}_{k=0}^{18}$ are polynomials chosen for physical interpretability.
On the lattice, these polynomials are \emph{uniquely determined} by the 19 numerical values in each row of $\Mmat$.
However, their \emph{continuous polynomial extensions} (i.e., their expressions as functions of continuous variables $X,Y,Z \in \Real$) are \textbf{not unique} due to the lattice reduction identities~\eqref{eq:lattice-reduction}.

We illustrate with the energy-squared moment $P_2$ (row index 2).
The numerical values of $\Mmat$ give:
\[
  \Mmat[2,:] = (12,\; -4,\; -4,\; -4,\; -4,\; -4,\; -4,\; 1,\; 1,\; 1,\; 1,\; 1,\; 1,\; 1,\; 1,\; 1,\; 1,\; 1,\; 1).
\]

Two valid continuous polynomial extensions that reproduce these values on the lattice are:

\medskip
\noindent\textbf{(A) Dubois continuous form} (keeping $r^4 = (X^2+Y^2+Z^2)^2$ as a true degree-4 polynomial):
\begin{equation}\label{eq:P2-continuous}
  P_2^{\text{cont}}(X,Y,Z) = \tfrac{21}{2}\,r^4 - \tfrac{53}{2}\,r^2 + 12.
\end{equation}

\noindent\textbf{(B) Lattice-reduced form} (applying $X^4 \to X^2$, etc.):
\begin{equation}\label{eq:P2-reduced}
  P_2^{\text{red}}(X,Y,Z) = -16\,r^2 + 21\,(X^2 Y^2 + X^2 Z^2 + Y^2 Z^2) + 12.
\end{equation}

\begin{proposition}\label{prop:equiv-on-lattice}
$P_2^{\text{cont}}(\vect{e}_i) = P_2^{\text{red}}(\vect{e}_i)$ for all $i = 0,\ldots,18$.
\end{proposition}
\begin{proof}
On the lattice, $X^4 = X^2$, so $r^4 = (X^2+Y^2+Z^2)^2 = X^2 + Y^2 + Z^2 + 2(X^2 Y^2 + X^2 Z^2 + Y^2 Z^2)$ (using $X^4 = X^2$).
Substituting into~\eqref{eq:P2-continuous}:
$\frac{21}{2}[r^2 + 2(X^2 Y^2 + X^2 Z^2 + Y^2 Z^2)] - \frac{53}{2}\,r^2 + 12 = -16\,r^2 + 21(X^2 Y^2 + X^2 Z^2 + Y^2 Z^2) + 12$.
This is exactly~\eqref{eq:P2-reduced}.
\end{proof}

% ============================================================
\section{Complete D3Q19 Polynomial Basis: Two Versions}\label{sec:polys}
% ============================================================

We now list all 19 polynomials in both forms.
The rows are grouped by physical type.
Polynomials of degree $\leq 2$ have unique extensions (no ambiguity); the difference arises only in the 6 polynomials containing degree-3 or degree-4 terms.

\subsection{Polynomials with unique extension (13 rows)}

These polynomials have degree $\leq 3$ using only the 19 independent monomials, so both versions agree:

\begin{center}
\renewcommand{\arraystretch}{1.3}
\begin{tabular}{clll}
\toprule
$k$ & Label & Type & $P_k(X,Y,Z)$ \\
\midrule
0 & $\rho$ (density) & conserved & $1$ \\
1 & $e$ (energy) & ghost & $19\,r^2 - 30$ \\
3 & $j_x$ & conserved & $X$ \\
5 & $j_y$ & conserved & $Y$ \\
7 & $j_z$ & conserved & $Z$ \\
9 & $3p_{xx}$ (stress) & stress & $2X^2 - Y^2 - Z^2$ \\
11 & $p_{ww}$ (stress) & stress & $Y^2 - Z^2$ \\
13 & $p_{xy}$ (stress) & stress & $XY$ \\
14 & $p_{yz}$ (stress) & stress & $YZ$ \\
15 & $p_{xz}$ (stress) & stress & $XZ$ \\
16 & $m_x$ (3rd) & 3rd-order & $X(Y^2 - Z^2)$ \\
17 & $m_y$ (3rd) & 3rd-order & $Y(Z^2 - X^2)$ \\
18 & $m_z$ (3rd) & 3rd-order & $Z(X^2 - Y^2)$ \\
\bottomrule
\end{tabular}
\end{center}

Here $r^2 \equiv X^2 + Y^2 + Z^2$.
These are degree $\leq 3$ and involve only the 19 independent monomials~\eqref{eq:19-monomials}.

\subsection{Polynomials with ambiguous extension (6 rows)}

\renewcommand{\arraystretch}{1.4}
\begin{center}
\begin{tabular}{cp{0.42\textwidth}p{0.42\textwidth}}
\toprule
$k$ & \textbf{Dubois Continuous} $P_k^{\text{cont}}$ & \textbf{Lattice-Reduced} $P_k^{\text{red}}$ \\
\midrule
2 &
$\frac{21}{2}\,r^4 - \frac{53}{2}\,r^2 + 12$
&
$-16\,r^2 + 21(X^2 Y^2{+}X^2 Z^2{+}Y^2 Z^2) + 12$
\\[4pt]
4 &
$5X\,r^2 - 9X$
\newline {\small (keeps $X^3$)}
&
$-4X + 5XY^2 + 5XZ^2$
\newline {\small ($X^3 \to X$)}
\\[4pt]
6 &
$5Y\,r^2 - 9Y$
&
$-4Y + 5YX^2 + 5YZ^2$
\\[4pt]
8 &
$5Z\,r^2 - 9Z$
&
$-4Z + 5ZX^2 + 5ZY^2$
\\[4pt]
10 &
$(2X^2{-}Y^2{-}Z^2)(3r^2{-}5)$
\newline {\small (keeps $X^4$, $X^2 Y^2$ as deg-4)}
&
$-4X^2{+}2Y^2{+}2Z^2 + 3(X^2 Y^2{+}X^2 Z^2) - 6Y^2 Z^2$
\\[4pt]
12 &
$(Y^2{-}Z^2)(3r^2{-}5)$
&
$-2Y^2{+}2Z^2 + 3X^2 Y^2 - 3X^2 Z^2$
\\
\bottomrule
\end{tabular}
\end{center}

\begin{remark}
On the D3Q19 lattice points, every $P_k^{\text{cont}}(\vect{e}_i) = P_k^{\text{red}}(\vect{e}_i)$ exactly, by the same argument as Proposition~\ref{prop:equiv-on-lattice}.
The two forms are indistinguishable when evaluated at $\vect{e}_i$.
They differ only when evaluated at shifted points $\vect{e}_i - \tilde{\vect{u}}$, which is precisely what happens in the construction of $\Tmat(\tilde{\vect{u}})$.
\end{remark}

% ============================================================
\section{Construction of \texorpdfstring{$\Tmat(\tilde{\vect{u}})$}{T(u)}}\label{sec:construction}
% ============================================================

\subsection{General definition (Dubois et al.\ 2015)}

Given a polynomial basis $\{P_k\}$, define the shifted moment matrix:
\begin{equation}\label{eq:M-shifted}
  M_{kj}(\tilde{\vect{u}}) = P_k(\vect{e}_j - \tilde{\vect{u}}), \qquad 0 \leq k,j \leq 18.
\end{equation}
The shift operator is:
\begin{equation}\label{eq:T-general}
  \Tmat(\tilde{\vect{u}}) = \Mmat(\tilde{\vect{u}})\,\Mmat^{-1}(0).
\end{equation}
This matrix converts ``fixed'' raw moments $\vect{m}(0) = \Mmat(0)\,\vect{f}$ into ``shifted'' moments $\vect{m}(\tilde{\vect{u}}) = \Tmat(\tilde{\vect{u}})\,\vect{m}(0)$.

\subsection{Why polynomial choice matters}

The shifted point is $\vect{e}_j - \tilde{\vect{u}}$, which is \emph{not} a lattice point (since $\tilde{\vect{u}} \in \Real^3$ is continuous).
At this point, $X^4 \neq X^2$ in general.
Therefore, $P_k^{\text{cont}}(\vect{e}_j - \tilde{\vect{u}}) \neq P_k^{\text{red}}(\vect{e}_j - \tilde{\vect{u}})$ for rows $k \in \{2,4,6,8,10,12\}$, leading to two different shift operators:
\begin{equation}
  \Tmat^{\text{cont}}(\tilde{\vect{u}}) \neq \Tmat^{\text{red}}(\tilde{\vect{u}}) \qquad \text{for } \tilde{\vect{u}} \neq \vect{0}.
\end{equation}

\subsection{Block-triangular structure}

Both versions of $\Tmat(\tilde{\vect{u}})$ share the same block-triangular structure (by degree ordering).
The d'Humi\`eres polynomials have degrees:
\begin{center}
\begin{tabular}{lll}
\toprule
Degree & Row indices & Count \\
\midrule
0 & $\{0\}$ & 1 \\
1 & $\{3,5,7\}$ & 3 \\
2 & $\{1,9,11,13,14,15\}$ & 6 \\
3 & $\{4,6,8,16,17,18\}$ & 6 \\
4 & $\{2,10,12\}$ & 3 \\
\bottomrule
\end{tabular}
\end{center}

Since shifting $P_k(\vect{X} - \tilde{\vect{u}})$ introduces terms of degree $\leq \deg(P_k)$, the matrix $\Tmat$ is block lower-triangular with respect to this degree ordering, with identity blocks on the diagonal.

% ============================================================
\section{The Lattice-Reduced Shift Operator: Closed Form}\label{sec:closed-form}
% ============================================================

We now derive the complete closed-form expressions.
Denote $\tilde{\vect{u}} = (\tu_x, \tu_y, \tu_z)$ and abbreviate $u_x^2 = \tu_x^2$, $u^2 = \tu_x^2 + \tu_y^2 + \tu_z^2$, etc.

The shift operator acts on the raw moment vector $\vect{m} = (m_0, m_1, \ldots, m_{18})^{\!\top}$ to produce the central moment vector $\vect{k} = (k_0, k_1, \ldots, k_{18})^{\!\top}$:
\begin{equation}
  \vect{k} = \Tmat(\tilde{\vect{u}})\,\vect{m}.
\end{equation}

\subsection{Conserved moments (degree 0 and 1, rows 0,3,5,7)}

These are independent of the polynomial extension choice:
\begin{align}
  k_0 &= m_0, \label{eq:k0}\\
  k_3 &= -\tu_x\,m_0 + m_3, \label{eq:k3}\\
  k_5 &= -\tu_y\,m_0 + m_5, \label{eq:k5}\\
  k_7 &= -\tu_z\,m_0 + m_7. \label{eq:k7}
\end{align}

\begin{remark}
For $\vect{f} = \feq$, we have $m_0 = \rho$ and $m_{3,5,7} = \rho\,\tu_{x,y,z}$, giving $k_0 = \rho$ and $k_{3,5,7} = 0$, confirming that central moments of the equilibrium have zero momentum.
\end{remark}

\subsection{Stress moments (degree 2, rows 9,11,13,14,15)}

Also independent of the polynomial extension (degree $\leq 2$, no ambiguity):
\begin{align}
  k_9  &= (2\tu_x^2 - \tu_y^2 - \tu_z^2)\,m_0 - 4\tu_x\,m_3 + 2\tu_y\,m_5 + 2\tu_z\,m_7 + m_9, \label{eq:k9}\\
  k_{11} &= (\tu_y^2 - \tu_z^2)\,m_0 - 2\tu_y\,m_5 + 2\tu_z\,m_7 + m_{11}, \label{eq:k11}\\
  k_{13} &= \tu_x\tu_y\,m_0 - \tu_y\,m_3 - \tu_x\,m_5 + m_{13}, \label{eq:k13}\\
  k_{14} &= \tu_y\tu_z\,m_0 - \tu_z\,m_5 - \tu_y\,m_7 + m_{14}, \label{eq:k14}\\
  k_{15} &= \tu_x\tu_z\,m_0 - \tu_z\,m_3 - \tu_x\,m_7 + m_{15}. \label{eq:k15}
\end{align}

\subsection{Energy moment (degree 2, row 1)}

\begin{equation}\label{eq:k1}
  k_1 = 19\,u^2\,m_0 + m_1 - 38\tu_x\,m_3 - 38\tu_y\,m_5 - 38\tu_z\,m_7.
\end{equation}

\subsection{Third-order moments (rows 16,17,18)}

These polynomials have degree 3 and only involve monomials from the independent set~\eqref{eq:19-monomials}, so no ambiguity:
\begin{align}
  k_{16} &= \tu_x(\tu_z^2 - \tu_y^2)\,m_0 + (\tu_y^2 - \tu_z^2)\,m_3 + 2\tu_x\tu_y\,m_5 - 2\tu_x\tu_z\,m_7 \nonumber\\
  &\quad - \tu_x\,m_{11} - 2\tu_y\,m_{13} + 2\tu_z\,m_{15} + m_{16}, \label{eq:k16}\\[6pt]
  k_{17} &= \tu_y(\tu_x^2 - \tu_z^2)\,m_0 - 2\tu_x\tu_y\,m_3 + (\tu_z^2 - \tu_x^2)\,m_5 + 2\tu_y\tu_z\,m_7 \nonumber\\
  &\quad + \tfrac{1}{2}\tu_y\,m_9 + \tfrac{1}{2}\tu_y\,m_{11} + 2\tu_x\,m_{13} - 2\tu_z\,m_{14} + m_{17}, \label{eq:k17}\\[6pt]
  k_{18} &= \tu_z(\tu_y^2 - \tu_x^2)\,m_0 + 2\tu_x\tu_z\,m_3 - 2\tu_y\tu_z\,m_5 + (\tu_x^2 - \tu_y^2)\,m_7 \nonumber\\
  &\quad - \tfrac{1}{2}\tu_z\,m_9 + \tfrac{1}{2}\tu_z\,m_{11} + 2\tu_y\,m_{14} - 2\tu_x\,m_{15} + m_{18}. \label{eq:k18}
\end{align}

\subsection{Energy-flux moments (degree 3/4, rows 4,6,8) --- \textbf{version-dependent}}

Using the \textbf{lattice-reduced} polynomials $P_4^{\text{red}} = -4X + 5XY^2 + 5XZ^2$ (with $X^3 \to X$):

\begin{align}
  k_4 &= \bigl(-5\tu_x\tu_y^2 - 5\tu_x\tu_z^2 - \tfrac{24}{19}\tu_x\bigr)\,m_0
       - \tfrac{10}{57}\tu_x\,m_1 \nonumber\\
  &\quad + 5(\tu_y^2 + \tu_z^2)\,m_3 + m_4
       + 10\tu_x\tu_y\,m_5 + 10\tu_x\tu_z\,m_7 \nonumber\\
  &\quad + \tfrac{5}{3}\tu_x\,m_9
       - 10\tu_y\,m_{13} - 10\tu_z\,m_{15}, \label{eq:k4}
\end{align}
and analogously for $k_6$ (cyclic: $x \to y$, $m_9 \to -\frac{5}{6}m_9 + \frac{5}{2}m_{11}$) and $k_8$ ($x \to z$, $m_9 \to -\frac{5}{6}m_9 - \frac{5}{2}m_{11}$).

\subsection{Ghost/energy-squared moments (degree 4, rows 2,10,12) --- \textbf{version-dependent}}

Using the \textbf{lattice-reduced} polynomial $P_2^{\text{red}}$:
\begin{align}
  k_2 &= \bigl[21(\tu_x^2\tu_y^2 + \tu_x^2\tu_z^2 + \tu_y^2\tu_z^2) + \tfrac{116}{19}\,u^2\bigr]\,m_0
       + \tfrac{14}{19}\,u^2\,m_1 + m_2 \nonumber\\
  &\quad + \bigl(-42\tu_x\tu_y^2 - 42\tu_x\tu_z^2 - \tfrac{8}{5}\tu_x\bigr)\,m_3
       - \tfrac{42}{5}\tu_x\,m_4 \nonumber\\
  &\quad + \bigl(-42\tu_x^2\tu_y - 42\tu_y\tu_z^2 - \tfrac{8}{5}\tu_y\bigr)\,m_5
       - \tfrac{42}{5}\tu_y\,m_6 \nonumber\\
  &\quad + \bigl(-42\tu_x^2\tu_z - 42\tu_y^2\tu_z - \tfrac{8}{5}\tu_z\bigr)\,m_7
       - \tfrac{42}{5}\tu_z\,m_8 \nonumber\\
  &\quad + (-7\tu_x^2 + \tfrac{7}{2}\tu_y^2 + \tfrac{7}{2}\tu_z^2)\,m_9
       + (-\tfrac{21}{2}\tu_y^2 + \tfrac{21}{2}\tu_z^2)\,m_{11} \nonumber\\
  &\quad + 84\tu_x\tu_y\,m_{13} + 84\tu_y\tu_z\,m_{14} + 84\tu_x\tu_z\,m_{15}. \label{eq:k2}
\end{align}

Rows 10 and 12 follow the same structure with coefficients derived from $P_{10}^{\text{red}}$ and $P_{12}^{\text{red}}$ (see Appendix~\ref{app:full}).

% ============================================================
\section{Inverse Shift Operator \texorpdfstring{$\Tmat^{-1}(\tilde{\vect{u}})$}{T inverse}}\label{sec:inverse}
% ============================================================

\subsection{Lattice-reduced version: exact inverse via sign flip}

\begin{theorem}[Group property]\label{thm:group}
For the lattice-reduced polynomial basis, the shift operator satisfies:
\begin{equation}\label{eq:group}
  \Tmat^{\text{red}}(\vect{u})\,\Tmat^{\text{red}}(\vect{v}) = \Tmat^{\text{red}}(\vect{u}+\vect{v}), \qquad \forall\,\vect{u},\vect{v} \in \Real^3.
\end{equation}
In particular, $\Tmat^{\text{red}}(-\vect{u}) = [\Tmat^{\text{red}}(\vect{u})]^{-1}$.
\end{theorem}

\begin{proof}
The lattice-reduced polynomials use only the 19 independent monomials~\eqref{eq:19-monomials}.
Denote $\vect{X} = (X,Y,Z)$.
The key observation is that the map $\vect{X} \mapsto \vect{X} - \vect{u}$ acts as a polynomial automorphism on the ring $\Real[X,Y,Z] / \mathcal{I}$, where $\mathcal{I}$ is the ideal generated by $\{X^4 - X^2, Y^4 - Y^2, Z^4 - Z^2, X^3 - X, Y^3 - Y, Z^3 - Z\}$.

In this quotient ring, the shift $(X,Y,Z) \to (X-u_x, Y-u_y, Z-u_z)$ followed by $(X,Y,Z) \to (X-v_x, Y-v_y, Z-v_z)$ is the same as a single shift $(X,Y,Z) \to (X-u_x-v_x, Y-u_y-v_y, Z-u_z-v_z)$.
Since $\Mmat^{-1}(0)$ maps into the quotient ring, the composition $\Tmat(\vect{u})\,\Tmat(\vect{v})$ equals $\Tmat(\vect{u}+\vect{v})$.

This has been verified symbolically and numerically to machine precision.
\end{proof}

\begin{corollary}
The inverse shift operator $\Tmat^{-1}(\vect{u})$ is obtained by replacing $\vect{u} \to -\vect{u}$ in all expressions.
Concretely, for the lattice-reduced version, the inverse $\vect{m} = \Tmat^{-1}(\vect{u})\,\vect{k}$ is obtained by:
\begin{enumerate}[nosep]
  \item Flipping the sign of all odd-power terms in $\tu_x, \tu_y, \tu_z$.
  \item Exchanging the roles of $\vect{m}$ and $\vect{k}$.
\end{enumerate}
For example, from~\eqref{eq:k3}: $k_3 = -\tu_x\,m_0 + m_3$ becomes $m_3 = +\tu_x\,k_0 + k_3$.
\end{corollary}

\subsection{Dubois continuous version: approximate inverse}

\begin{proposition}\label{prop:dubois-no-group}
For the Dubois continuous polynomial basis, the group property fails:
\begin{equation}
  \Tmat^{\text{cont}}(\vect{u})\,\Tmat^{\text{cont}}(-\vect{u}) \neq \Imat.
\end{equation}
The error is $O(|\vect{u}|^2)$ and is concentrated in rows $\{2, 4, 6, 8, 10, 12\}$.
\end{proposition}

\begin{proof}
The continuous polynomial $r^4 = (X^2+Y^2+Z^2)^2$ has degree 4 in the unreduced ring.
At shifted points $\vect{e}_i - \vect{u}$, expanding $(e_{ix}-u_x)^4$ introduces the monomial $u_x^4$, which is \emph{not} in the lattice-reduced ring (it would be $u_x^2$ under lattice reduction, but $u_x$ is continuous, not a lattice variable).

More precisely, $\Mmat^{-1}(0)$ was constructed from the 19 lattice-independent monomials.
The continuous polynomial $P_k^{\text{cont}}(\vect{e}_j - \vect{u})$ introduces ``ghost'' monomials $e_{jx}^4 - 4e_{jx}^3 u_x + \cdots$ that project onto the 19-dimensional lattice space through aliasing: $e_{jx}^4 \to e_{jx}^2$.
This aliasing breaks the automorphism property.

Numerical verification: at $\vect{u} = (0.12, -0.07, 0.05)$,
\[
  \|\Tmat^{\text{cont}}(\vect{u})\,\Tmat^{\text{cont}}(-\vect{u}) - \Imat\|_F = 0.345.
\]
\end{proof}

% ============================================================
\section{Comparison of the Two Versions}\label{sec:comparison}
% ============================================================

\subsection{What is identical}

\begin{corollary}[Dubois et al.\ 2015, Corollary 2.1]\label{cor:viscosity}
The second-order equivalent partial differential equations (and hence the Navier--Stokes viscosity) are independent of the polynomial extension choice.
Specifically, the stress-mode central moments (rows 9, 11, 13, 14, 15) are \emph{identical} between both versions:
\begin{equation}
  [\Tmat^{\text{cont}}(\vect{u})]_{k,:} = [\Tmat^{\text{red}}(\vect{u})]_{k,:}, \qquad k \in \{0,1,3,5,7,9,11,13,14,15,16,17,18\}.
\end{equation}
\end{corollary}

This follows because these rows correspond to polynomials of degree $\leq 3$ that involve only the 19 independent monomials --- no degree-4 ambiguity exists.

\subsection{What differs}

The 6 rows $\{2, 4, 6, 8, 10, 12\}$ differ.
These are the ``starred'' ghost modes (energy-squared, energy-flux, energy$\times$stress).
The difference in $\Tmat$ is an additive correction $\Delta\Tmat$:
\begin{equation}
  \Tmat^{\text{cont}} = \Tmat^{\text{red}} + \Delta\Tmat, \qquad \text{where } \Delta\Tmat \text{ has nonzero entries only in rows } \{2,4,6,8,10,12\}.
\end{equation}

\subsection{Correction structure}

The correction $\Delta\Tmat$ can be expressed in terms of the ``correction polynomials'':
\begin{equation}\label{eq:delta-gamma}
  \delta(\alpha) = \alpha^2(\alpha^2 - 1), \qquad \gamma(\alpha) = \alpha(\alpha^2 - 1).
\end{equation}
These vanish on the lattice ($\alpha \in \{-1,0,1\} \Rightarrow \delta = \gamma = 0$) but are nonzero at shifted points.

The correction terms are proportional to $\delta(\tu_x), \gamma(\tu_x)$, etc., and thus scale as $O(\tu^2)$ for small $\tu$.
In particular, at $\tu = 0$, $\Delta\Tmat = 0$ (both versions degenerate to the identity).

\subsection{Summary table}

\begin{center}
\renewcommand{\arraystretch}{1.3}
\begin{tabular}{lcc}
\toprule
Property & Lattice-Reduced & Dubois Continuous \\
\midrule
Group: $\Tmat(\vect{u})\Tmat(-\vect{u}) = \Imat$ & \textcolor{green!60!black}{\textbf{Exact}} & \textcolor{red}{\textbf{Fails}} ($O(u^2)$) \\
Morphism: $\Tmat(\vect{u}+\vect{v}) = \Tmat(\vect{u})\Tmat(\vect{v})$ & \textcolor{green!60!black}{\textbf{Exact}} & \textcolor{red}{\textbf{Fails}} ($O(u^2)$) \\
Inverse: $\Tmat^{-1} = \Tmat(-\vect{u})$ & \textcolor{green!60!black}{\textbf{Yes}} & \textcolor{red}{\textbf{No}} (requires matrix inv.) \\
Stress CMs & Identical & Identical \\
N-S viscosity & Identical & Identical \\
Ghost/energy CMs & Version A & Version B $\neq$ A \\
Stability & $\checkmark$ & Possibly different \\
\bottomrule
\end{tabular}
\end{center}

% ============================================================
\section{Physical Interpretation}
% ============================================================

The MRT-CM collision~\eqref{eq:collision} relaxes the non-equilibrium central moments with rates $s_k$.
The conserved central moments ($k_0, k_3, k_5, k_7$) have $s_k = 0$ (not relaxed).
The stress central moments ($k_9, k_{11}, k_{13}, k_{14}, k_{15}$) are relaxed with rate $s_{\nu}$ that determines the kinematic viscosity:
\begin{equation}
  \nu = \frac{1}{3}\Bigl(\frac{1}{s_\nu} - \frac{1}{2}\Bigr)\,\Delta x^2 / \Delta t.
\end{equation}

Since the stress CMs are identical between both versions (Corollary~\ref{cor:viscosity}), the physical viscosity is unchanged regardless of which polynomial extension is used.

The ghost modes (rows 2, 4, 6, 8, 10, 12) affect only higher-order error terms and stability margins.
The lattice-reduced version, with its exact algebraic structure, is expected to provide more consistent stability behavior, particularly at higher Mach numbers where $|\Delta\Tmat|$ becomes significant.

% ============================================================
\section{Implementation Notes}\label{sec:implementation}
% ============================================================

\subsection{Efficient CUDA implementation}

The shift operator does \textbf{not} require a $19\times 19$ matrix--vector multiply.
Both $\Tmat(\vect{u})\,\vect{m}$ and $\Tmat^{-1}(\vect{u})\,\vect{k}$ can be evaluated as hand-coded arithmetic expressions (Section~\ref{sec:closed-form}), requiring approximately:

\begin{center}
\begin{tabular}{lcc}
\toprule
 & Multiplications & Additions \\
\midrule
$\Tmat\,\vect{m}$ (forward) & $\sim 120$ & $\sim 100$ \\
$\Tmat^{-1}\,\vect{k}$ (inverse) & $\sim 120$ & $\sim 100$ \\
\bottomrule
\end{tabular}
\end{center}

This is far cheaper than a full $19\times 19$ matrix multiply ($361 \times 2 = 722$ FLOPs).

\subsection{Compile-time switch}

In our GILBM code, a compile-time macro selects the version:
\begin{verbatim}
  #define USE_DUBOIS_CONTINUOUS  0  // 0 = lattice-reduced, 1 = Dubois continuous
\end{verbatim}

The Dubois version is implemented as an additive correction to the lattice-reduced version:
\begin{verbatim}
  raw_to_central_dH_dubois(m, ux, uy, uz, k):
      raw_to_central_dH(m, ux, uy, uz, k);     // lattice-reduced base
      k[2]  += correction_row2(m, ux, uy, uz);  // additive DeltaT
      k[4]  += correction_row4(m, ux, uy, uz);
      ...
      k[12] += correction_row12(m, ux, uy, uz);
\end{verbatim}

This structure minimizes code duplication and allows direct comparison.

% ============================================================
\section{Conclusion}
% ============================================================

We have presented a complete derivation of the central moment shift operator $\Tmat(\tilde{\vect{u}})$ for D3Q19 MRT lattice Boltzmann methods, filling a gap in the literature where Dubois et al.\ (2015) only provided explicit D2Q9 expressions.

The key insight is that the d'Humi\`eres D3Q19 polynomial basis has a \textbf{non-unique continuous extension} due to the lattice identity $x^4 = x^2$ (for $x \in \{-1,0,1\}$).
We identified two natural choices:
\begin{enumerate}[nosep]
  \item The ``Dubois continuous'' form, which preserves $r^4 = (r^2)^2$ as a true degree-4 polynomial.
  \item The ``lattice-reduced'' form, which uses only 19 independent monomials by applying $x^4 \to x^2$, $x^3 \to x$.
\end{enumerate}

The lattice-reduced form is algebraically superior: it satisfies the group property $\Tmat(\vect{u})\,\Tmat(-\vect{u}) = \Imat$ exactly, and provides an explicit inverse $\Tmat^{-1}(\vect{u}) = \Tmat(-\vect{u})$ without matrix inversion.
Both forms yield identical stress-mode central moments and hence the same Navier--Stokes viscosity.

% ============================================================
\appendix
\section{Full Expressions for All Remaining Rows}\label{app:full}
% ============================================================

\subsection{Row 10: $3\pi_{xx}$ (energy $\times$ stress)}

\begin{align}
  k_{10} &= \bigl[3(\tu_x^2\tu_y^2 + \tu_x^2\tu_z^2) - 6\tu_y^2\tu_z^2 - \tfrac{16}{19}\tu_x^2 + \tfrac{8}{19}\tu_y^2 + \tfrac{8}{19}\tu_z^2\bigr]\,m_0 \nonumber\\
  &\quad + \tfrac{1}{19}(2\tu_x^2 - \tu_y^2 - \tu_z^2)\,m_1 \nonumber\\
  &\quad + (-6\tu_x\tu_y^2 - 6\tu_x\tu_z^2 + \tfrac{16}{5}\tu_x)\,m_3 - \tfrac{6}{5}\tu_x\,m_4 \nonumber\\
  &\quad + (-6\tu_x^2\tu_y + 12\tu_y\tu_z^2 - \tfrac{8}{5}\tu_y)\,m_5 + \tfrac{3}{5}\tu_y\,m_6 \nonumber\\
  &\quad + (-6\tu_x^2\tu_z + 12\tu_y^2\tu_z - \tfrac{8}{5}\tu_z)\,m_7 + \tfrac{3}{5}\tu_z\,m_8 \nonumber\\
  &\quad + (-\tu_x^2 + 2\tu_y^2 + 2\tu_z^2)\,m_9 + m_{10} + (3\tu_y^2 - 3\tu_z^2)\,m_{11} \nonumber\\
  &\quad + 12\tu_x\tu_y\,m_{13} - 24\tu_y\tu_z\,m_{14} + 12\tu_x\tu_z\,m_{15} \nonumber\\
  &\quad + 9\tu_y\,m_{17} - 9\tu_z\,m_{18}. \label{eq:k10}
\end{align}

\subsection{Row 12: $\pi_{ww}$ (energy $\times$ stress)}

\begin{align}
  k_{12} &= \bigl[3\tu_x^2(\tu_y^2 - \tu_z^2) + \tfrac{8}{19}(\tu_z^2 - \tu_y^2)\bigr]\,m_0 + \tfrac{1}{19}(\tu_y^2 - \tu_z^2)\,m_1 \nonumber\\
  &\quad - 6\tu_x(\tu_y^2 - \tu_z^2)\,m_3 \nonumber\\
  &\quad + \tfrac{2}{5}\tu_y(4 - 15\tu_x^2)\,m_5 - \tfrac{3}{5}\tu_y\,m_6 \nonumber\\
  &\quad + \tfrac{2}{5}\tu_z(15\tu_x^2 - 4)\,m_7 + \tfrac{3}{5}\tu_z\,m_8 \nonumber\\
  &\quad + (\tu_y^2 - \tu_z^2)\,m_9 + 3\tu_x^2\,m_{11} + m_{12} \nonumber\\
  &\quad + 12\tu_x\tu_y\,m_{13} - 12\tu_x\tu_z\,m_{15} \nonumber\\
  &\quad - 6\tu_x\,m_{16} + 3\tu_y\,m_{17} + 3\tu_z\,m_{18}. \label{eq:k12}
\end{align}

\subsection{Row 6: $q_y$ (energy flux, $y$-direction)}

\begin{align}
  k_6 &= (-5\tu_x^2\tu_y - 5\tu_y\tu_z^2 - \tfrac{24}{19}\tu_y)\,m_0
       - \tfrac{10}{57}\tu_y\,m_1 \nonumber\\
  &\quad + 10\tu_x\tu_y\,m_3 + 5(\tu_x^2 + \tu_z^2)\,m_5 + m_6
       + 10\tu_y\tu_z\,m_7 \nonumber\\
  &\quad - \tfrac{5}{6}\tu_y\,m_9 + \tfrac{5}{2}\tu_y\,m_{11}
       - 10\tu_x\,m_{13} - 10\tu_z\,m_{14}. \label{eq:k6}
\end{align}

\subsection{Row 8: $q_z$ (energy flux, $z$-direction)}

\begin{align}
  k_8 &= (-5\tu_x^2\tu_z - 5\tu_y^2\tu_z - \tfrac{24}{19}\tu_z)\,m_0
       - \tfrac{10}{57}\tu_z\,m_1 \nonumber\\
  &\quad + 10\tu_x\tu_z\,m_3 + 10\tu_y\tu_z\,m_5
       + 5(\tu_x^2 + \tu_y^2)\,m_7 + m_8 \nonumber\\
  &\quad - \tfrac{5}{6}\tu_z\,m_9 - \tfrac{5}{2}\tu_z\,m_{11}
       - 10\tu_y\,m_{14} - 10\tu_x\,m_{15}. \label{eq:k8}
\end{align}

% ============================================================
\section{Dubois Correction Terms \texorpdfstring{$\Delta\Tmat$}{Delta T}}\label{app:correction}
% ============================================================

The additive correction $\Delta\Tmat$ satisfies $\Tmat^{\text{cont}}\,\vect{m} = \Tmat^{\text{red}}\,\vect{m} + \Delta\Tmat\,\vect{m}$.
Only rows $\{2,4,6,8,10,12\}$ are nonzero, and each row's correction depends only on $m_0, m_1, m_3, m_5, m_7, m_9, m_{11}$ (at most 7 moment entries):

\medskip
\noindent\textbf{Row 2:}
\begin{align}
  \Delta k_2 &= \bigl[\tfrac{21}{2}\tu_x^4 + \tfrac{861}{38}\tu_x^2 + \tfrac{21}{2}\tu_y^4 + \tfrac{861}{38}\tu_y^2 + \tfrac{21}{2}\tu_z^4 + \tfrac{861}{38}\tu_z^2\bigr]\,m_0 \nonumber\\
  &\quad + \tfrac{21}{19}(\tu_x^2 + \tu_y^2 + \tu_z^2)\,m_1 \nonumber\\
  &\quad + (-42\tu_x^3 - 21\tu_x)\,m_3 + (-42\tu_y^3 - 21\tu_y)\,m_5 + (-42\tu_z^3 - 21\tu_z)\,m_7 \nonumber\\
  &\quad + (21\tu_x^2 - \tfrac{21}{2}\tu_y^2 - \tfrac{21}{2}\tu_z^2)\,m_9 + (\tfrac{63}{2}\tu_y^2 - \tfrac{63}{2}\tu_z^2)\,m_{11}. \label{eq:delta2}
\end{align}

\noindent\textbf{Row 4:}
\begin{equation}\label{eq:delta4}
  \Delta k_4 = (-5\tu_x^3 - \tfrac{55}{19}\tu_x)\,m_0 - \tfrac{5}{19}\tu_x\,m_1 + 15\tu_x^2\,m_3 - 5\tu_x\,m_9.
\end{equation}

\noindent\textbf{Row 6:}
\begin{equation}\label{eq:delta6}
  \Delta k_6 = (-5\tu_y^3 - \tfrac{55}{19}\tu_y)\,m_0 - \tfrac{5}{19}\tu_y\,m_1 + 15\tu_y^2\,m_5 + \tfrac{5}{2}\tu_y\,m_9 - \tfrac{15}{2}\tu_y\,m_{11}.
\end{equation}

\noindent\textbf{Row 8:}
\begin{equation}\label{eq:delta8}
  \Delta k_8 = (-5\tu_z^3 - \tfrac{55}{19}\tu_z)\,m_0 - \tfrac{5}{19}\tu_z\,m_1 + 15\tu_z^2\,m_7 + \tfrac{5}{2}\tu_z\,m_9 + \tfrac{15}{2}\tu_z\,m_{11}.
\end{equation}

\noindent\textbf{Row 10:}
\begin{align}
  \Delta k_{10} &= (6\tu_x^4 + \tfrac{246}{19}\tu_x^2 - 3\tu_y^4 - \tfrac{123}{19}\tu_y^2 - 3\tu_z^4 - \tfrac{123}{19}\tu_z^2)\,m_0 \nonumber\\
  &\quad + \tfrac{1}{19}(12\tu_x^2 - 6\tu_y^2 - 6\tu_z^2)\,m_1 \nonumber\\
  &\quad + (-24\tu_x^3 - 12\tu_x)\,m_3 + (12\tu_y^3 + 6\tu_y)\,m_5 + (12\tu_z^3 + 6\tu_z)\,m_7 \nonumber\\
  &\quad + (12\tu_x^2 + 3\tu_y^2 + 3\tu_z^2)\,m_9 + (-9\tu_y^2 + 9\tu_z^2)\,m_{11}. \label{eq:delta10}
\end{align}

\noindent\textbf{Row 12:}
\begin{align}
  \Delta k_{12} &= (3\tu_y^4 + \tfrac{123}{19}\tu_y^2 - 3\tu_z^4 - \tfrac{123}{19}\tu_z^2)\,m_0 + \tfrac{1}{19}(6\tu_y^2 - 6\tu_z^2)\,m_1 \nonumber\\
  &\quad + (-12\tu_y^3 - 6\tu_y)\,m_5 + (12\tu_z^3 + 6\tu_z)\,m_7 \nonumber\\
  &\quad + (-3\tu_y^2 + 3\tu_z^2)\,m_9 + (9\tu_y^2 + 9\tu_z^2)\,m_{11}. \label{eq:delta12}
\end{align}

% ============================================================
\section*{References}
% ============================================================

\begin{enumerate}[label={[\arabic*]},nosep]
  \item F.~Dubois, T.~Fevrier, and B.~Graille, ``Lattice Boltzmann schemes with relative velocities,'' \textit{Commun.\ Comput.\ Phys.}, vol.~17, no.~4, pp.~1088--1112, 2015. arXiv:1312.3297.
  \item D.~d'Humi\`eres, I.~Ginzburg, M.~Krafczyk, P.~Lallemand, and L.-S.~Luo, ``Multiple-relaxation-time lattice Boltzmann models in three dimensions,'' \textit{Phil.\ Trans.\ R.\ Soc.\ Lond.\ A}, vol.~360, pp.~437--451, 2002.
  \item M.~Geier, A.~Greiner, and J.~G.~Korvink, ``Cascaded digital lattice Boltzmann automata for high Reynolds number flow,'' \textit{Phys.\ Rev.\ E}, vol.~73, p.~066705, 2006.
  \item T.~Imamura, K.~Suzuki, T.~Nakamura, and M.~Yoshida, ``Acceleration of steady-state lattice Boltzmann simulations on non-uniform mesh using local time step method,'' \textit{J.~Comput.\ Phys.}, vol.~202, pp.~645--663, 2005.
\end{enumerate}

\end{document}
